\section{Quantitative Results}

This section evaluates the model quantitatively and links the equilibrium mechanisms to simulated outcomes. The numerical solution uses the fully nested global-game formulation developed in the code. For each overvaluation state, the algorithm jointly solves for the run cutoff, aggregate foreign-currency demand, the allocation between cash and stablecoins, and the endogenous precision of the public signal.

\subsection{Implementation}

The model is solved on a grid of overvaluation states. For each state, we integrate over public-signal realizations and compute the exact posterior beliefs of the marginal runner. This yields a fully nested fixed point in which stablecoin usage and signal precision are determined jointly. The resulting equilibrium objects are total foreign-currency demand, stablecoin usage, crisis incidence, and welfare.

\subsection{Foreign-Currency Demand}

The first result is that foreign-currency demand increases with overvaluation. At low levels of overvaluation, this increase is driven primarily by the hedge motive. As overvaluation rises further, the run component becomes quantitatively important, increasing the sensitivity of aggregate demand to fundamentals.

Stablecoin demand also increases with overvaluation. As foreign-currency demand rises, stablecoins account for an increasing share of total demand because they remain the lower-cost instrument over the relevant region. This result is central for the paper's empirical interpretation: stablecoin activity should be strongest precisely when misalignment and parallel market pressure are most severe.

\subsection{Crisis Outcomes}

The second result is that stablecoins increase crisis risk. In the cash-only benchmark, crises arise when hedge demand and run demand together exceed effective defense capacity. Introducing stablecoins shifts the economy toward a higher incidence of crises. This occurs through two channels. First, stablecoins reduce the effective cost of acquiring foreign currency and therefore raise baseline demand. Second, stablecoins improve the precision of public information, increasing coordination on the run margin.

The interaction of these two channels is important. A reduction in access costs alone would increase demand, but would not generate the same amplification of coordinated runs. Conversely, better public information without lower access costs would sharpen coordination but would not generate the same increase in aggregate positions. The quantitative results show that both channels matter jointly.

\subsection{Welfare}

The third result is that the welfare effects of stablecoins are non-monotonic. In low and moderate overvaluation states, stablecoins improve welfare because the cost of acquiring foreign currency falls and households can better smooth access to tradable goods. In high-overvaluation states, however, welfare declines relative to the cash-only benchmark. In this region, the increase in crisis frequency and the greater severity of crisis-state losses dominate the access benefit.

This welfare reversal is the central normative implication of the model. Stablecoins are privately useful but may become socially harmful when fundamentals are weak and coordination effects are strong.

\subsection{Mechanism Decomposition}

The quantitative exercise also clarifies the distinct roles of the cost and information channels. Reducing the cost advantage of stablecoins weakens the increase in foreign-currency demand and attenuates welfare gains in low-overvaluation states. Weakening the information channel reduces coordination, narrows the crisis region, and weakens the welfare losses at high overvaluation. These counterfactuals show that the paper's main results are not driven by a single force but by the interaction of access and coordination.

\subsection{Summary}

Overall, the quantitative results confirm the theory. Stablecoins increase foreign-currency demand, increase crisis risk, and generate a welfare tradeoff that depends on the underlying state of the economy. The model therefore provides a coherent explanation for why stablecoins may simultaneously expand financial access and increase macro-financial fragility in parallel exchange-rate environments.
